\chapter{양자 요동과 구조의 기원}

\lectureoverview{마지막 강의에서는 인플레이션으로 매끄러워진 우주가 어떻게 은하가 자라나는 씨앗인 미세한 불균일성을 얻게 되는지 묻는다. 답은 여러 단계로 구성된다. 먼저 보편 상수로 플랑크 척도를 정하고, 익숙한 입자물리학과 양자중력 사이에 놓인 거대한 간극을 드러낸다. 이어 우주 마이크로파 배경은 목표값을 제시한다. 여러 척도에 걸쳐 퍼져 있는 온도와 밀도의 변동은 대략 \(10^{-5}\)이다. 영점 운동은 불가피한 요동을 제공하고, 감쇠는 진동이 어떻게 얼어붙을 수 있는지 설명한다. 팽창 우주에서의 파동 방정식은 인플라톤의 모든 푸리에 모드를 바로 그런 감쇠 진동자로 바꾼다. 물리적 파장이 허블 척도에 이르면 모드는 얼어붙는다. 인플레이션이 끝날 때 이 동결된 변동 때문에 이웃한 영역들이 서로 조금씩 다른 시점에 재가열되고, 그 결과 미시적인 양자 잡음이 원시 밀도 분포로 바뀐다.}

\section{어떤 상수가 보편적인가?}

우리는 기본 입자에 관한 매우 정밀한 양자 이론을 갖고 있다. 이 이론의 매개변수에는 쿼크와 렙톤의 질량, 결합 상수, 혼합 자료, 전하, 그리고 입자 목록 자체가 포함된다. 계산 능력이 무한하다면 이 이론은 가속기 실험뿐 아니라 원자, 분자, 화학까지 설명할 수 있을 것이다. 원칙적으로 그 결과는 생물학과 사고의 기반이 되는 물리적 과정까지 이어진다. 그렇다고 이런 계를 계산하기 쉽다는 뜻은 아니다. 표준 모형은 보통 물질의 미시적 토대가 되도록 의도된 이론이라는 뜻이다. 중력은 고전적 수준에서 일반 상대성 이론으로 별도로 기술된다. 두 이론 모두 강력하지만, 저에너지 물리학에 나타나는 모든 수치가 필연적인 이유를 설명해 주지는 않는다.

양성자와 전자의 질량비는 물리적으로 잘 정의된 수치지만, 특정한 두 입자에 결부되어 있다. 양성자를 중성자로, 전자를 뮤온으로 바꾸면 다른 비가 나온다. 이런 매개변수에는 우리가 아직 모르는 심오한 물리가 담겨 있을 수 있지만, 그 자체로 보편성을 드러내지는 않는다. 다음 세 가지 차원 상수는 다르다.
\[
c,\qquad \hbar,\qquad G.
\]
이 상수들을 정의하는 법칙에는 \emph{모든}이라는 말이 거듭 등장한다. 모든 물질과 정보는 한계 속도 \(c\)를 따른다. 모든 양자계는 \(\hbar\)가 정하는 불확정성 관계를 만족한다. 모든 형태의 에너지는 동일한 뉴턴 상수 \(G\)에 따라 중력을 만든다. 이처럼 폭넓게 적용된다는 점에서 이 논의에서 이 상수들을 기본 상수로 다룬다.

양자역학의 온전한 의미가 알려지기 전부터 플랑크는 이 조합의 중요성을 알아차렸다. 그의 상수는 처음에 흑체 복사를 기술하는 데 등장했다. 그때도 차원에 관한 사실은 분명했다. \(c\), \(\hbar\), \(G\)로부터 특정한 원자나 입자에 의존하지 않는 길이, 시간, 질량의 단위를 구성할 수 있었다.

\section{플랑크 척도의 구성}

길이, 시간, 질량의 차원을 각각 \(L\), \(T\), \(M\)이라 하자. 빛의 속도 차원은 다음과 같다.
\[
[c]=LT^{-1}.
\]
불확정성 관계 \(\Delta x\,\Delta p\gtrsim\hbar\)와 \(p=mv\)을 함께 쓰면 다음을 얻는다.
\[
[\hbar]=ML^2T^{-1}.
\]
마지막으로 뉴턴 방정식
\[
F=\frac{Gm_1m_2}{r^2},
\qquad
F=ma
\]
에서 다음을 얻는다.
\[
[G]=L^3M^{-1}T^{-2}.
\]
이 유도가 중요한 이유는 차원 분석이 추측이 아니기 때문이다. 불필요한 차원이 소거되도록 각 지수를 정하면 된다.

다음 형태의 길이를 구해 보자.
\begin{equation}
\ell=c^\alpha\hbar^\beta G^\gamma.
\label{eq:c10-length-ansatz}
\end{equation}
질량 지수는 0이어야 하므로
\[
\beta-\gamma=0.
\]
시간 지수는 0이어야 하고, 길이 지수는 1이어야 한다.
\[
-\alpha-\beta-2\gamma=0,
\qquad
\alpha+2\beta+3\gamma=1.
\]
\(\gamma=\beta\)을 이용하면 다음과 같다.
\[
\alpha+3\beta=0,
\qquad
\alpha+5\beta=1.
\]
두 식의 차를 취하면 \(2\beta=1\)이고, 따라서
\[
\beta=\gamma=\frac12,
\qquad
\alpha=-\frac32.
\]
세 보편 상수로 만들 수 있는 유일한 길이는 다음과 같다.
\begin{equation}
\ell_{\mathrm P}
=\sqrt{\frac{\hbar G}{c^3}}
\approx 1.6\times10^{-35}\,\mathrm m.
\label{eq:c10-planck-length}
\end{equation}
강의에서는 차수만 따져 \(10^{-35}\,\mathrm m\)로 나타내지만, 차원 유도는 상수들의 조합 자체를 정해 준다.

이에 대응하는 시간은 빛이 이 길이를 가로지르는 데 걸리는 시간이다.
\begin{equation}
t_{\mathrm P}
=\frac{\ell_{\mathrm P}}{c}
=\sqrt{\frac{\hbar G}{c^5}}
\approx 5.4\times10^{-44}\,\mathrm s.
\label{eq:c10-planck-time}
\end{equation}
차수만 따지면 \(10^{-43}\,\mathrm s\)이다. 플랑크 길이와 보편적인 속도를 알면 이를 독립적으로 따로 구성할 필요는 없다.

질량은 차원 관계 \(Et\sim\hbar\)와 \(E=mc^2\)을 이용하면 된다. \(t=t_{\mathrm P}\)을 대입하면
\begin{equation}
m_{\mathrm P}
=\frac{\hbar}{c^2t_{\mathrm P}}
=\sqrt{\frac{\hbar c}{G}}
\approx 2.2\times10^{-8}\,\mathrm{kg}.
\label{eq:c10-planck-mass}
\end{equation}
입자의 질량 기준으로 보면 놀라울 만큼 거시적인 값으로, 대략 고운 먼지 한 점의 질량에 해당한다. 정지 에너지는 \(10^9\,\mathrm J\) 정도로, 연료 탱크에 든 화학 에너지와 비슷한 규모다. 연소는 연료의 정지 에너지 가운데 극히 일부만 방출하므로 자동차는 연료를 몇 탱크만 써도 수백 마일을 달릴 수 있다. 플랑크 질량 하나를 전부 에너지로 바꾸는 것은 전혀 다른 일이다. 작은 질량도 \(c^2\)를 곱하면 막대한 에너지를 나타낼 수 있다.

가속기 물리학과의 간극은 크다. 강의 당시 고에너지 충돌에서 탐사한 가장 짧은 척도의 대표값인 \(10^{-18}\,\mathrm m\)의 거리는 여전히 플랑크 길이의 약 \(10^{17}\)배다. 따라서 전약 척도와 플랑크 척도 사이의 미지 영역은 광대하다. 핵물리학이 좋은 비유가 된다. 멀리서 보면 원자핵은 기본 입자처럼 보이지만, 더 짧은 거리에서 보면 양성자와 중성자로 분해되고, 이들은 다시 쿼크와 글루온으로 분해된다. 저에너지 입자는 원자핵에 비하면 단순할 수 있지만, 플랑크 척도에서 무엇이 존재하든 그것에 견주면 여전히 복잡한 창발적 대상일 수 있다. 이 계층 구조를 고려하면 표준 모형에 매개변수가 많은 점은 덜 놀랍지만, 관측된 값과 미세 조정의 이유는 여전히 설명되지 않는다.

\begin{classroomqa}
\audiencequestion{힉스 입자는 예상보다 가벼운가, 무거운가? 질량이 왜 수수께끼인가?}
\lecturerresponse{비교 기준은 두 가지다. 플랑크 질량을 기준으로 보면 \(125\,\mathrm{GeV}\)인 힉스 입자는 매우 가볍다. 보호 구조가 없는 이론에서는 양자 보정으로 스칼라 질량이 훨씬 높은 척도에 민감해지는 경향이 있으므로, 관측된 값이 나오려면 매개변수들이 아슬아슬하게 균형을 이뤄야 하는 것처럼 보인다. 반면 당시 논의된 몇 가지 단순한 초대칭 이론의 예측과 비교하면 관측된 힉스 입자를 다소 무겁다고 부를 수도 있다. 겉보기의 모순은 서로 다른 두 기준 척도와 비교한 데서 비롯되며, 근본적인 계층성 문제는 여전히 남는다. \lecturetimestamp{00:26:45--00:30:54}}
\end{classroomqa}

\begin{classroomqa}
\audiencequestion{힉스 입자가 입자에 질량을 부여하는 데 관여한다면, 암흑 물질 입자도 플랑크 척도에 가까운 질량을 가져야 하나?}
\lecturerresponse{아니다. 윔프는 \emph{약하게 상호작용하는 무거운 입자}를 뜻한다. 이 표현이 가리키는 것은 실험실에서 검출할 수 있을 가능성을 열어 주는 상호작용 척도이지, 플랑크 질량이 아니다. 통상적인 윔프 가정 아래에서는 양성자 질량보다 훨씬 큰 질량도 가능하지만, 여전히 \(m_{\mathrm P}\)보다 엄청나게 작다. 보통 물질과의 상호작용을 임의로 더 약하게 허용하면 우주론적으로 가능한 질량 범위는 훨씬 넓어지지만, 그만큼 직접 검출은 더 어려워진다. 양성자와 중성자는 유용한 주의 사례이기도 하다. 둘의 질량 대부분은 힉스가 만들어 낸 쿼크 질량의 단순한 합이 아니라 강한 상호작용의 결합 에너지다. \lecturetimestamp{00:30:55--00:33:11}}
\end{classroomqa}

\begin{classroomqa}
\audiencequestion{플랑크 질량은 무엇을 뜻하는가? 비트당 근본적인 에너지가 있는가?}
\lecturerresponse{중력에 관한 유용한 해석은, 통상적인 준고전적 블랙홀로 남아 있는 한 블랙홀의 크기를 매개변수에 따라 플랑크 척도보다 훨씬 작게 만들 수 없다는 것이다. 그러한 블랙홀의 질량은 \(m_{\mathrm P}\) 정도이고, 지평선 넓이는 플랑크 면적 하나 정도이며, 엔트로피도 단위 하나 정도다. 이처럼 제한된 블랙홀 맥락에서는 정보 단위 하나 정도에 플랑크 에너지를 연관 지을 수 있다. 그러나 모든 장치에 저장된 모든 비트에 플랑크 에너지가 든다는 보편적 주장은 아니다. 보통 입자와 일반적인 기억 장치는 훨씬 가벼울 수 있다. \lecturetimestamp{00:33:11--00:35:10}}
\end{classroomqa}

이제 플랑크 단위를 쓰면 편리하다.
\begin{equation}
c=\hbar=G=1.
\label{eq:c10-planck-units}
\end{equation}
그러면 길이, 시간, 질량, 에너지를 ~\eqref{eq:c10-planck-length}--\eqref{eq:c10-planck-mass}의 방정식에 견준 순수한 수치로 나타낼 수 있다.

\section{거의 균일한 마이크로파 하늘}

우주론의 문제는 하늘 전체를 담은 마이크로파 지도로부터 시작된다. 윌킨슨 마이크로파 비등방성 탐사선인 WMAP은 원시 플라스마가 투명해진 시기 무렵에 방출된 복사를 측정했다. 그 뒤 광자는 대략 천 배 적색편이되어 현재 약 \(2.7\,\mathrm K\)의 거의 완벽한 흑체 복사를 이룬다. 바깥을 볼 때 우리는 우주를 현재의 한 시점에서 보는 것이 아니다. 거리가 멀수록 빛이 더 이른 시기에 방출되었으며, 마이크로파 배경은 광자가 우리를 향해 자유롭게 날아오기 시작한 먼 구면에서 온 빛이다.

\begin{classroomqa}
\audiencequestion{타원형 WMAP 이미지는 정확히 무엇을 나타내는가?}
\lecturerresponse{타원 안의 각 점은 하늘의 한 방향을 나타내며, 구면을 지도 투영법으로 펼쳐 놓은 것이다. 좌우 가장자리는 서로 이어지고, 위아래 부분은 은하의 서로 반대쪽 반구를 나타내며, 은하면은 가운데를 가로지른다. 우리 은하에서 나온 전경 마이크로파 방출은 모형화해 빼야 한다. 남는 것은 가까운 별들의 지도가 아니라 마지막 산란 구면에서 온 마이크로파 온도를 방향별로 나타낸 지도다. \lecturetimestamp{00:37:49--00:40:49}}
\end{classroomqa}

일차적으로 결과는 등방적이다.
\[
T(\hat{\mathbf n})\simeq 2.7\,\mathrm K
\]
모든 관측 방향 \(\hat{\mathbf n}\)에서 그렇다. 모든 지점을 중심으로 완벽한 등방성이 성립한다면 완벽한 균질성도 뜻하게 된다. 하지만 완벽하게 균질한 우주에서는 은하와 은하단, 그리고 지난 강의에서 다룬 초점면 같은 우주 거대 구조가 저절로 생기지 않는다. 진폭이 아무리 작더라도 어떤 비균질성은 존재했어야 한다.

방향에 따른 온도 편차를 다음과 같이 쓴다
\[
\delta T(\hat{\mathbf n})=T(\hat{\mathbf n})-\bar T.
\]
측정된 분수 변동의 크기는 다음 정도이다
\begin{equation}
\frac{\delta T}{T}\sim10^{-5}.
\label{eq:c10-cmb-anisotropy}
\end{equation}
이 값에 이르는 한 가지 방법은 동역학적이었다. 오늘날의 물질 분포를 과거로 거슬러 진화시켜 탈결합 시점에 어떤 초기 밀도 대비가 필요했는지 묻는 것이다. 다른 방법은 직접 측정하는 것이었다. 거의 균일한 마이크로파 지도를 다섯 자릿수 규모로 증폭해 그 얼룩무늬 자체를 측정했다. 두 방법은 이 논증에 필요한 수준에서 일치한다.

얼룩무늬는 여러 각도 척도에 걸쳐 나타난다. 큰 얼룩 안에 더 작은 얼룩이 있고, 한 영역을 확대하면 통계적으로 비슷한 질감이 보인다. \emph{척도 불변}이라는 말은 모든 척도가 문자 그대로 같다는 뜻이 아니다. 관련 범위에서 요동의 파워가 척도에 따라 아주 약하게만 변한다는 뜻이다. 이러한 거의 척도 불변적인 성질은 이론이 설명해야 할 증거의 일부이지, 단지 전체 진폭만을 가리키는 것이 아니다.

온도는 왜 물질을 추적해야 할까? 중력 퍼텐셜 우물을 빠져나오는 광자는 에너지를 잃는다. 따라서 밀도가 높은 영역에서 오는 복사는 밀도가 낮은 영역에서 오는 복사와 비교해 편이될 수 있다. 강의에서 다루는 차수 수준에서는 이로부터 다음 관계를 예상할 수 있다.
\begin{equation}
\frac{\delta\rho}{\rho}\sim10^{-5}
\label{eq:c10-density-target}
\end{equation}
분리 시점 무렵에 그렇다. 정밀한 우주 마이크로파 배경 계산에는 음향 동역학, 중력 퍼텐셜, 전달 함수가 포함되므로 방정식~\eqref{eq:c10-density-target}은 \(\delta T/T\)와 \(\delta\rho/\rho\) 사이의 정확한 등식이 아니라 척도를 가늠하기 위한 약식 표현이다.

\begin{classroomqa}
\audiencequestion{얼룩무늬는 마지막 산란 때 생긴 보통의 미시적 잡음일 수 있는가?}
\lecturerresponse{관측된 각도 척도에서는 그럴 수 없다. 분리 시점에 해당 패치들은 재결합 플라스마의 원자 규모 교란이 아니라 엄청나게 큰 물리적 영역이었다. 하지만 시간을 훨씬 더 거슬러 올라가면 물리적 파장은 더 작아진다. 인플레이션이 여기에 필요한 확대를 제공한다. 처음에는 양자적이고 미시적인 교란이었던 요동을 천문학적 척도까지 늘릴 수 있다. \lecturetimestamp{00:46:42--00:48:02}}
\end{classroomqa}

이 메커니즘은 세 가지 물리 개념으로 이해할 수 있다. 조화 진동자의 영점 운동, 진동자의 감쇠, 그리고 팽창 우주에서 스칼라장의 파동 방정식이다. 차례로 살펴보자.

\section{진동자의 영점 운동}

단위 질량을 가진 조화 진동자의 경우
\begin{equation}
E=\frac12\dot x^{\,2}+\frac12\omega^2x^2,
\qquad
\ddot x+\omega^2x=0.
\label{eq:c10-undamped}
\end{equation}
운동 에너지와 퍼텐셜 에너지는 끊임없이 서로 바뀐다. 한 번 진동하는 동안 평균하면 각각 전체 에너지의 절반을 이룬다.

양자역학에서 에너지 준위는 다음과 같다.
\begin{equation}
E_n=\left(n+\frac12\right)\hbar\omega.
\label{eq:c10-oscillator-levels}
\end{equation}
바닥 상태에서는 \(x=0\)과 \(\dot x=0\)을 동시에 가질 수 없다. 그렇게 되면 위치와 운동량이 동시에 정확하게 정해지기 때문이다. 바닥 상태에 남는 최소 에너지
\[
E_0=\frac12\hbar\omega,
\]
를 영점 에너지라 한다. \emph{진동}이라는 말은 통계적인 의미로 이해해야 한다. 정상 바닥 상태는 고전적인 사인파 궤적을 그리지 않지만 위치와 운동량의 분산은 0이 아니다.

평균적으로 바닥 상태 에너지의 절반은 퍼텐셜 항에 들어 있다.
\[
\frac12\omega^2\langle x^2\rangle
=\frac14\hbar\omega.
\]
따라서
\begin{equation}
\langle x^2\rangle=\frac{\hbar}{2\omega}
\sim\frac{\hbar}{\omega}.
\label{eq:c10-zero-point-position}
\end{equation}
2와 같은 계수는 무시하면 요동의 크기를 복소수 표기로 다음처럼 나타낼 수 있다.
\[
x(t)\sim\sqrt{\frac{\hbar}{\omega}}\,e^{i\omega t}.
\]
이를 미분하면
\begin{equation}
\langle\dot x^{\,2}\rangle\sim\hbar\omega,
\label{eq:c10-zero-point-velocity}
\end{equation}
를 얻으며, 이는 운동 에너지와 퍼텐셜 에너지의 평균이 같다는 점과 일치한다. 냉각하면 \(n>0\)인 들뜬 상태는 없앨 수 있지만, 바닥 상태의 이러한 분산까지 없앨 수는 없다.

\section{감쇠, 임계 상태, 동결}

진동자를 점성이 있는 매질에 넣어 보자. 꿀 속에서 움직이는 용수철과 질량은 다음 방정식을 따른다.
\begin{equation}
\ddot x+\gamma\dot x+\omega^2x=0,
\label{eq:c10-damped}
\end{equation}
여기서 \(\gamma>0\)는 항력을 나타낸다. \(x=e^{\alpha t}\)를 대입하면
\begin{equation}
\alpha^2+\gamma\alpha+\omega^2=0,
\end{equation}
를 얻고, 따라서
\begin{equation}
\alpha_\pm
=\frac{-\gamma\pm\sqrt{\gamma^2-4\omega^2}}{2}.
\label{eq:c10-damped-roots}
\end{equation}
이다. 두 근은 2계 미분 방정식에 필요한 서로 독립적인 두 해를 제공한다.

\begin{classroomqa}
\audiencequestion{방정식~\eqref{eq:c10-damped}의 항력 항에 마이너스 부호가 있어야 하지 않나요?}
\lecturerresponse{항력 자체는 \(F_{\rm drag}=-\gamma\dot x\)이다. 뉴턴 방정식은 \(\ddot x=-\gamma\dot x-\omega^2x\)이다. 두 힘 항을 모두 왼쪽으로 옮기면 방정식~\eqref{eq:c10-damped}처럼 플러스 부호가 된다. 따라서 미분 방정식의 부호와 물리적 힘의 방향은 서로 일치한다. \lecturetimestamp{00:59:38--01:00:10}}
\end{classroomqa}

\(\gamma^2<4\omega^2\)이면 근은 복소수다. 다음과 같이 정의하자.
\[
\Omega=\sqrt{\omega^2-\frac{\gamma^2}{4}}.
\]
그러면
\begin{equation}
x(t)=e^{-\gamma t/2}
\left(A\cos\Omega t+B\sin\Omega t\right).
\label{eq:c10-underdamped}
\end{equation}
이다. 진동자는 여전히 평형점을 반복해서 지나지만, 진폭을 감싸는 포락선은 줄어든다. 이것이 저감쇠 영역이다.

\(\gamma^2>4\omega^2\)이면 방정식~\eqref{eq:c10-damped-roots}의 두 근은 모두 실수이며 음수다.
\begin{equation}
x(t)=A e^{\alpha_+t}+B e^{\alpha_-t}.
\label{eq:c10-overdamped}
\end{equation}
지속적인 진동은 없고 두 성분 모두 감쇠한다. 이것이 과감쇠 영역이다. \begin{equation}
\gamma=2\omega
\label{eq:c10-critical}
\end{equation}일 때는 두 근이 일치하며, 두 거동 사이의 경계인 임계 감쇠가 된다.

특히 중요한 극한이 하나 있다. 복원력이 없으면 \(\omega=0\)이고,
\[
\alpha(\alpha+\gamma)=0,
\]
이므로
\begin{equation}
x(t)=c_0+c_1e^{-\gamma t}.
\label{eq:c10-no-restoring}
\end{equation}
이다. 속도는 사라지지만 좌표가 0으로 돌아갈 필요는 없다. 초기 조건이 정하는 상수 \(c_0\)에서 멈춘 채로 남는다. 용수철 없이 꿀 속에서 밀린 돌은 항력 때문에 운동을 멈춘 위치에 그대로 멈춘다.

이제 진동수가 천천히 변한다고 하자. 먼저 \(\omega\gg\gamma\)인 상태에서 시작하면 운동은 진동한다. 감쇠 척도는 거의 일정하게 유지하면서 \(\omega(t)\)를 낮춘다. 진동은 느려지고 \(\gamma=2\omega\)을 지나 과감쇠 상태가 되며, \(\omega\to0\)일 때 복원력이 없는 형태에 가까워진다. 또는 복원력에 비해 감쇠를 키워도 된다. 우주론에 적용할 때 결정적인 변화는 허블 감쇠가 거의 일정한 가운데 유효 진동수가 낮아지는 것이다. 동결 후의 값은 이 경계만으로 알 수 없다. 이전의 전체적인 시간 변화와 초기 양자 상태로부터 정해진다.

\section{팽창 우주의 스칼라 파동}

인플라톤 \(\phi\)으로 돌아가자. 퍼텐셜은 높고 넓은 평탄부 뒤에 낮은 최솟값으로 이어지는 가파른 내리막이 있다고 가정한다. 평탄부에서 중요한 근사는 다음과 같다.
\begin{equation}
V'(\phi)\simeq0,
\label{eq:c10-flat-potential}
\end{equation}
\(V(\phi)\simeq0\)가 아니다. 에너지 밀도가 인플레이션을 일으킬 만큼 큰 채로 유지될 수 있는 동안 기울기는 작다. 이 형태의 세부적인 기원은 주어지지 않는다. 이는 모형이 현상론적으로 요구하는 조건이다.

\begin{figure}[htbp]\centering
\includegraphics[width=0.88\textwidth]{lecture_10_figure_02.png}\caption{왼쪽에는 인플라톤의 평탄부와 절벽이 보이고, 오른쪽에서는 평탄한 시공간의 파동 방정식이 시작된다. 칠판에 적힌 내용으로는 근사 조건 \(V'(\phi)=0\)과 이계 미분 항의 시작 부분을 확인할 수 있지만, 강의자가 완성되지 않은 방정식을 가려 놓았다. 방정식~\eqref{eq:c10-flat-wave}은 구두 유도를 바탕으로 재구성했다.}\label{fig:c10-potential-wave}\lectureframe{01:17:17.865}\end{figure}

평탄한 시공간에서 거의 자유로운 스칼라장의 경우,
\begin{equation}
\frac{\partial^2\phi}{\partial t^2}
-\nabla^2\phi=0.
\label{eq:c10-flat-wave}
\end{equation}
퍼텐셜이 평탄하지 않으면 오른쪽에 \(-V'(\phi)\)가 추가된다. 평탄부 근사에서는 이 항을 생략한다. 장이 완전히 균질하다면 공간 미분은 0이 된다. 팽창 우주에서는 균질한 장에 대해 남는 방정식이
\[
\ddot\phi+3H\dot\phi=0,
\]
이 되며, 이는 지난 강의의 허블 마찰 방정식이다.

\begin{figure}[htbp]\centering
\includegraphics[width=0.88\textwidth]{lecture_10_figure_03.png}\caption{왼쪽에는 평탄부가 표시되고, 오른쪽에는 파동 방정식이 단계적으로 나타난다. 위쪽의 평탄한 시공간 연산자는 읽을 수 있지만, 아래쪽 식은 \(+3H\)까지 적힌 채 미완성이다. 따라서 전체 감쇠 항은 당시의 설명을 바탕으로 채웠다.}\label{fig:c10-hubble-board}\lectureframe{01:18:57.554}\end{figure}

\begin{figure}[htbp]\centering
\begin{coursefitfigure}\begin{tikzpicture}[x=1cm,y=0.9cm,>=Latex]
  \draw[->] (-0.2,0) -- (9.3,0) node[right] {\(\phi\)};
  \draw[->] (0,-0.2) -- (0,4.6) node[above] {\(V(\phi)\)};
  \draw[very thick,blue!55!black,smooth]
    plot coordinates {(0.25,3.9) (1.4,3.88) (2.8,3.8) (4.2,3.6) (5.3,3.25) (6.0,2.6) (6.5,1.45) (6.9,0.45) (7.5,0.28) (8.5,0.7)};
  \draw[dashed] (6.45,0) -- (6.45,2.0);
  \draw[->,thick,red!65!black] (2.0,4.25) -- (3.1,3.82);
  \node[anchor=south west] at (0.8,4.0) {slow-roll plateau: \(V'\simeq0\)};
  \node[anchor=west] at (6.55,2.05) {end of inflation};
  \node[anchor=north] at (7.45,0.22) {reheating region};
\end{tikzpicture}\end{coursefitfigure}\caption{논의 전반에 사용된 퍼텐셜을 깨끗하게 재구성한 그림이다. 완만한 기울기에서는 느린 구름이 가능하고, \(V\)의 값이 큰 동안 \(H\)은 거의 일정하게 유지된다. 가파른 내리막에서 인플레이션이 끝나고 저장된 에너지가 방출된다.}\label{fig:c10-potential-reconstruction}\end{figure}

비균질한 장에서는 좌표 거리와 물리적 거리를 구분해야 한다. 무차원 공변 좌표 \(\mathbf x\)를 택하면 미소 고유 거리는 다음과 같다.
\[
d\boldsymbol\ell=a(t)\,d\mathbf x.
\]
시간을 고정하면
\[
\nabla_{\rm proper}=\frac{1}{a}\nabla_{\mathbf x},
\qquad
\nabla_{\rm proper}^2=\frac{1}{a^2}\nabla_{\mathbf x}^2.
\]
스칼라 방정식은 따라서 다음과 같다.
\begin{equation}
\ddot\phi+3H\dot\phi-\frac{1}{a^2}\nabla^2\phi=0,
\label{eq:c10-frw-wave}
\end{equation}
여기서 \(\nabla\)는 공변 좌표에 대한 미분을 나타낸다. 대신 무차원 척도 인자를 택하면 공변 좌표가 길이 차원을 갖는다. 물리는 달라지지 않는다.

\begin{classroomqa}
\audiencequestion{방정식~\eqref{eq:c10-frw-wave}의 공간 항은 양수여야 하지 않나요?}
\lecturerresponse{아니다. 평탄한 시공간의 파동 연산자에는 이미 상대 부호 \(\partial_t^2-\nabla^2\)가 들어 있다. 팽창은 공간 미분에 인자 \(a^{-2}\)를 곱하지만 부호를 뒤집지는 않는다. 푸리에 치환을 하면 두 마이너스 부호가 서로 상쇄되어 모드 방정식에는 예상대로 양의 복원 항 \(+k^2\phi_k/a^2\)이 생긴다. \lecturetimestamp{01:22:18--01:23:06}}
\end{classroomqa}

\section{각 푸리에 모드는 감쇠 진동자다}

공간 파동 하나씩 살펴보자.
\begin{equation}
\phi(\mathbf x,t)
=\phi_{\mathbf k}(t)e^{i\mathbf k\cdot\mathbf x}.
\label{eq:c10-mode-ansatz}
\end{equation}
복소 지수는 사인과 코사인을 간결하게 나타내기 위한 표기다. 임의의 실수 분포는 \(\mathbf k\)와 \(-\mathbf k\) 사이의 적절한 실수성 조건을 만족하는 모드들을 중첩해 구성한다.

시간 미분은 \(\phi_{\mathbf k}(t)\)에만 작용하고,
\[
\nabla^2e^{i\mathbf k\cdot\mathbf x}
=-k^2e^{i\mathbf k\cdot\mathbf x}.
\]
이를 방정식~\eqref{eq:c10-frw-wave}에 대입하면
\begin{equation}
\ddot\phi_{\mathbf k}
+3H\dot\phi_{\mathbf k}
+\frac{k^2}{a^2}\phi_{\mathbf k}=0.
\label{eq:c10-mode-equation}
\end{equation}
을 얻는다. 따라서 모든 파수는 다음과 같은 감쇠 진동자로 거동한다.
\begin{equation}
\gamma=3H,
\qquad
\omega_k(t)=\frac{k}{a(t)}.
\label{eq:c10-mode-identification}
\end{equation}
물리적 파수 \(k/a\)는 우주 팽창에 따라 적색편이된다. 그 결과 유효 복원력 \(k^2/a^2\)은 점차 0으로 줄어든다.

짧은 물리적 파장에서 시작하는 모드는 \(\omega_k\gg H\)를 만족하며 진동한다. 파장이 늘어나고 진동수는 낮아지며, 결국 허블 감쇠가 우세해진다. 그러면 0으로 되돌아가는 대신 진동하지 않는 값에 가까워진다. 서로 다른 \(\mathbf k\) 모드들은 각기 다른 시점에 경계를 넘는다. 더 큰 \(k\)는 \(a\)가 고정되어 있을 때 더 짧은 파장에 해당하므로, 물리적 진동수가 허블 척도까지 낮아지려면 더 오래 팽창해야 한다.

\section{지평선 이탈과 동결된 진공 요동}

진동자의 임계 조건을 차수 수준의 추정에 적용하면
\[
\gamma=2\omega_k,
\]
방정식~\eqref{eq:c10-mode-identification}으로부터
\begin{equation}
3H=\frac{2k}{a},
\qquad
a_{\rm cross}=\frac{2k}{3H}.
\label{eq:c10-critical-mode}
\end{equation}
를 얻는다. 수치 계수 자체가 핵심은 아니다. 물리적 파수 \(k/a\)가 \(H\) 정도가 될 때 경계 전이가 일어난다.

다음과 같이 두자.
\[
\lambda_{\rm com}=\frac{2\pi}{k},
\qquad
\lambda_{\rm phys}=a\lambda_{\rm com}
=\frac{2\pi a}{k}.
\]
차원을 확인하면 \(a\)의 위치를 정할 수 있다. \(\mathbf x\)가 무차원이면 \(k\)과 \(\mathbf k\cdot\mathbf x\)도 모두 무차원이고, \(a\)와 \(\lambda_{\rm phys}\)은 길이 차원을 갖는다. 따라서 방정식~\eqref{eq:c10-critical-mode}은 다음을 뜻한다.
\begin{equation}
\lambda_{\rm phys}\sim H^{-1},
\label{eq:c10-horizon-crossing}
\end{equation}
이는 \(2\), \(3\), \(2\pi\)과 같은 계수를 무시한 관계다.

허블 법칙 \(v=Hd\)에 따르면 거리 \(H^{-1}\)에 대응하는 후퇴 속도는 \(c=1\) 정도다. 인플레이션 중에는 \(H\)이 크고 거의 일정하므로, 이 허블 길이는 인과 척도 역할을 하면서도 미시적이다. 이 사건을 \emph{지평선 이탈}이라고 부르는 것은 관례적인 표현이다. 물질적 물체가 물리적 벽을 통과하는 것은 아니다. 모드의 물리적 파장이 거의 일정한 허블 척도보다 커지고, 그 동역학은 빠른 진동에서 거의 동결된 진폭으로 바뀐다.

\begin{figure}[htbp]\centering
\begin{coursefitfigure}\begin{tikzpicture}[x=1cm,y=0.85cm,>=Latex]
  \draw[->] (0,0) -- (9.2,0) node[right] {cosmic time};
  \draw[->] (0,0) -- (0,5.0) node[above] {physical length};
  \draw[dashed,thick] (0.25,2.75) -- (8.75,2.75)
    node[right] {\(H^{-1}\)};
  \draw[very thick,blue!55!black]
    (0.35,0.55) .. controls (3.3,0.75) and (5.3,1.55) .. (8.5,4.55);
  \node[blue!55!black,anchor=west] at (6.65,4.55) {\(\lambda_{\rm phys}=2\pi a/k\)};
  \draw[densely dotted] (6.05,0) -- (6.05,2.75);
  \fill (6.05,2.75) circle (2pt);
  \node[anchor=north] at (6.05,0) {horizon exit};
  \node[align=center] at (2.6,1.7) {oscillating\\sub-Hubble mode};
\end{tikzpicture}\end{coursefitfigure}\caption{거의 지수적인 팽창 중에는 허블 길이가 거의 일정하게 유지되는 동안 물리적 파장은 늘어난다. 둘이 만나는 시점에 대략적인 계수 차이를 제외하고 빠른 진동에서 동결 모드로 전이가 일어난다.}\label{fig:c10-horizon-reconstruction}\end{figure}

진공 요동이 있기 때문에 인플레이션이 모든 파동을 완전히 지워 버리지는 않는다. 파장이 짧은 모든 모드는 영점 운동으로 시작한다. 팽창에 따라 한 모드의 파장이 허블 척도까지 늘어나 얼어붙는 동안, 더 짧은 진공 모드들은 계속 진동한다. 이어 이 모드들도 늘어나 얼어붙는다. 이 과정이 반복되면서 다양한 파장과 방향에 걸친 동결된 성분을 포함하는 장의 분포가 만들어진다. 그 결과는 확률적이지만 무구조한 것은 아니다. 그 통계적 성질은 양자 상태와 인플레이션 배경으로부터 정해진다.

\begin{figure}[htbp]\centering
\includegraphics[width=0.88\textwidth]{lecture_10_figure_04.png}\caption{칠판 전체에 길고 짧은 장의 파동이 중첩되어 있고, \(\gamma\)와 \(3H\)라는 표기 옆에는 작은 감쇠 그림이 있다. 이 장면은 모드에 관한 정성적 그림을 보여 주며, 정확한 함수를 읽어 내는 데 쓰이지 않는다.}\label{fig:c10-mode-board}\lectureframe{01:45:28.476}\end{figure}

\begin{classroomqa}
\audiencequestion{모드는 한 지점에서 얼어붙는가, 파장에서 얼어붙는가? 그리고 \(H\) 자체가 지수적으로 증가하는가?}
\lecturerresponse{파장이 허블 척도에 이르면 얼어붙는다. 거의 평탄한 평탄부에서는 \(V\)와 그에 따른 \(H\)가 거의 일정하다. 지수적으로 커지는 것은 척도 인자 \(a(t)\)이며, 따라서 물리적 파장 \(a/k\)도 지수적으로 증가한다. 각 공변 모드는 저마다의 시점에 조건 \(k/a\sim H\)에 이른다. \lecturetimestamp{01:45:06--01:45:39}}
\end{classroomqa}

\begin{classroomqa}
\audiencequestion{전자기파인가요?}
\lecturerresponse{아니다. 인플라톤 스칼라장의 요동이다. 전자기 진공 요동은 익숙한 비유를 제공하지만, 방정식~\eqref{eq:c10-mode-equation}의 변수는 \(\phi\)다. \lecturetimestamp{01:45:39--01:46:04}}
\end{classroomqa}

\begin{classroomqa}
\audiencequestion{파장들이 연속적으로 분포하나요?}
\lecturerresponse{무한한 공간이라는 이상화에서는 그렇다. 장은 푸리에 적분으로 나타낸다.
\begin{equation}
\phi(\mathbf x,t)
=\int\frac{d^3k}{(2\pi)^3}\,
\phi_{\mathbf k}(t)e^{i\mathbf k\cdot\mathbf x}.
\label{eq:c10-fourier-integral}
\end{equation}
모든 파수에서 같은 적색편이와 동결 메커니즘이 일어난다. \lecturetimestamp{01:46:04--01:46:34}}
\end{classroomqa}

\begin{classroomqa}
\audiencequestion{우주가 유한하다면 모드는 이산적인가요?}
\lecturerresponse{그렇다. 유한한 공간 부피에 경계 조건을 부과하면 허용되는 파수는 양자화된다. 다루는 척도보다 훨씬 큰 부피에서는 간격이 매우 촘촘해져 그 합을 방정식~\eqref{eq:c10-fourier-integral}의 적분으로 잘 근사할 수 있다. 연속적인 그림은 이상화일 뿐, 국소적인 메커니즘 자체가 다른 것은 아니다. \lecturetimestamp{01:46:34--01:47:01}}
\end{classroomqa}

이제 두 운동이 함께 일어난다. 공간 평균장 \(\bar\phi(t)\)는 퍼텐셜을 따라 고전적으로 천천히 구른다. 그 위에 작은 양자 요동
\[
\phi(\mathbf x,t)=\bar\phi(t)+\delta\phi(\mathbf x,t),
\]
이 얹혀 있으며, 이는 동결된 모드들을 중첩해 만든 것이다. 요동 자체의 에너지는 관측된 물질 밀도 대비를 만들어 내기에는 너무 작다. 중요한 점은 이 요동이 인플레이션이 끝나는 시점을 교란한다는 것이다.

완전한 단일장 양자 계산을 하면 이 진동자 논의가 더 정밀해진다. 가벼운 스칼라장은 파수의 로그 간격마다 다음 정도의 대표적 요동을 얻는다.
\begin{equation}
\delta\phi_{\rm exit}\sim\frac{H}{2\pi}.
\label{eq:c10-field-amplitude}
\end{equation}
정규화 값을 유도하려면 인플레이션 배경에서 장의 모드를 양자화해야 한다. 진동자 그림은 이미 핵심 메커니즘을 보여 주었다. 바닥 상태에서도 진공 요동이 존재하며, 팽창은 이를 고전적으로 보이는 동결 분포로 바꾼다.

\section{절벽이 장의 잡음을 밀도 잡음으로 바꾼다}

먼저 \(\delta\phi\)를 억제하자. 그러면 어느 고정된 시각에 장은 어디서나 같은 값을 가지며, 공간 \(x\)을 위쪽에, 장의 값 \(\phi\)을 가로축에 놓으면 칠판에 수직선으로 나타난다. 시간이 흐르면서 이 선은 절벽 쪽으로 이동한다. 우주 전체가 함께 가장자리에 이르고, 인플라톤 퍼텐셜 에너지는 입자와 복사와 열로 바뀌며, 지수 팽창은 끝난다. 그 결과 생긴 에너지 밀도는 통상적인 팽창과 함께 희석되기 시작한다.

\begin{figure}[htbp]
\centering
\includegraphics[width=0.88\textwidth]{lecture_10_figure_05.png}
\caption{칠판에 그린 공간과 장의 값 사이의 도식. 수직축은 \(x\)이고, 가로 방향은 \(\phi\)이다. 희미한 수직에 가까운 곡선들은 연속적인 장의 구성을 나타내며, 오른쪽의 빗금 친 경계는 인플레이션이 끝나는 절벽이다.}
\label{fig:c10-cliff-board}
\lectureframe{01:50:54.502}
\end{figure}

\begin{figure}[htbp]
\centering
\begin{coursefitfigure}
\begin{tikzpicture}[x=1cm,y=0.82cm,>=Latex]
  \draw[->] (0,0) -- (9.2,0) node[right] {\(\phi\)};
  \draw[->] (0,0) -- (0,5.2) node[above] {space \(x\)};
  \draw[thick,gray!60] (8.2,0.15) rectangle (8.55,4.85);
  \foreach \y in {0.35,0.85,...,4.55} {
    \draw[gray!60] (8.2,\y) -- (8.55,{\y+0.25});
  }
  \node[anchor=south,rotate=90] at (8.38,2.5) {cliff};
  \draw[thick,densely dashed,gray!70] (3.2,0.2) -- (3.2,4.75);
  \draw[thick,blue!45!black,smooth]
    plot coordinates {(5.0,0.2) (5.12,0.75) (4.9,1.25) (5.18,1.8) (4.94,2.35) (5.2,2.9) (4.88,3.45) (5.12,4.05) (5.0,4.75)};
  \draw[very thick,red!65!black,smooth]
    plot coordinates {(7.15,0.2) (7.32,0.7) (7.0,1.2) (7.42,1.75) (7.1,2.3) (7.5,2.85) (7.02,3.35) (7.35,4.0) (7.16,4.75)};
  \draw[->,thick] (3.55,2.5) -- (4.55,2.5);
  \draw[->,thick] (5.55,2.5) -- (6.65,2.5);
  \node[align=center,anchor=south,font=\small] at (3.2,4.8) {uniform\\profile};
  \node[align=center,anchor=south,font=\small] at (5.0,4.8) {frozen\\ripple};
  \node[align=center,anchor=south,font=\small] at (7.2,4.8) {near\\reheating};
\end{tikzpicture}
\end{coursefitfigure}
\caption{공간과 장의 관계를 나타낸 그림을 정확히 재구성했다. 평균장은 절벽을 향해 움직이는 동안 공간적으로 고정된 잔물결은 지속되므로, 위치마다 인플레이션이 끝나는 시점이 조금씩 달라진다.}
\label{fig:c10-cliff-reconstruction}
\end{figure}

얼어붙은 변화를 되살려 보자. 이제 선은 약간 물결친다. 이웃한 위치마다 \(\phi\)의 값이 조금씩 다르다. 절벽을 향해 어깨를 나란히 맞추고 행진하려는 사람들의 대열을 떠올려 보자. 몇몇은 조금 앞서고, 몇몇은 조금 뒤처진다. 앞서 있는 사람들이 먼저 떨어진다.

이 시간 차이가 밀도 대비의 실제 근원이다. 먼저 가장자리에 이른 영역은 먼저 재가열된다. 그 영역의 퍼텐셜 에너지는 입자와 열로 전환되고, 우주가 인플레이션 이후 시대에 팽창하는 동안 그 에너지는 희석될 시간이 생긴다. 조금 뒤에 가장자리에 이른 인접 영역은 더 늦게 재가열된다. 따라서 두 영역을 비교하면, 나중에 가장자리에 이른 영역은 희석될 시간이 더 짧았다는 이유만으로 밀도가 더 높다. 그러므로 처음 일어난 통과는 상대적으로 밀도가 높은 영역들을 남긴다. 그 영역들이 팽창하는 동안, 더 짧은 파장의 잔물결은 다른 위치들을 가장자리 너머로 이끌어 더 큰 무늬 안에 더 작은 영역들을 만든다. 얼어붙은 모드의 전체 위계에 걸쳐 이 과정이 반복되면서 구조 안에 구조가 생긴다. 따라서 작은 장 요동이 중요한 것은 그 자체의 에너지 때문이 아니라, 재가열이 일어나는 국소적 시각을 바꾸기 때문이다.

시간 관계는 개략적으로 쓸 수 있다. 균질한 장이 속도 \(\dot{\bar\phi}\)로 최종값 \(\phi_{\rm end}\)에 이른다면, 작은 장 변위는 국소적인 종료 시각을 다음만큼 바꾼다.
\begin{equation}
\delta t_{\rm end}
\simeq-\frac{\delta\phi}{\dot{\bar\phi}}.
\label{eq:c10-time-delay}
\end{equation}
장이 느리게 굴러갈수록 같은 \(\delta\phi\)에 더 큰 지연이 대응한다. 그 추가 시간 동안 앞선 영역은 더 많이 팽창하고 희석된다. 1 정도의 전달 계수를 고려하면 곡률 섭동 또는 밀도 섭동은 다음 양으로 결정된다.
\begin{equation}
\mathcal R
\sim H\,\delta t_{\rm end}
\sim-\frac{H\,\delta\phi}{\dot{\bar\phi}}.
\label{eq:c10-curvature-estimate}
\end{equation}
이를 식~\eqref{eq:c10-field-amplitude}과 결합하면 익숙한 슬로 롤 스케일을 얻는다.
\begin{equation}
\mathcal P_{\mathcal R}^{1/2}
\sim\frac{H^2}{2\pi|\dot{\bar\phi}|},
\label{eq:c10-slow-roll-amplitude}
\end{equation}
정확한 규약과 이후 \(\delta\rho/\rho\)로의 전달을 다루려면 더 완전한 섭동 계산이 필요하다. 장이 느리게 굴러갈수록 이웃한 절벽 통과 사이의 시간은 길어지고, 그동안의 희석도 커진다. 따라서 같은 장 요동도 \(\lvert\dot{\bar\phi}\rvert\)가 작을수록 더 큰 밀도 흔적을 남긴다.

많은 모드가 서로 다른 시각에 허블 규모를 가로지르므로, 크기가 서로 다른 많은 영역이 조금씩 다른 시각에 재가열된다. 그 결과 밀도장은 여러 규모에 걸쳐 중첩된 과밀 영역과 저밀 영역을 포함한다. 밀도장의 통계적 성질은 인플레이션 퍼텐셜의 높이와 기울기와 곡률에 달려 있다. 적절한 느린 구름 모형은 자연스럽게 다음에 가까운 규모 불변 진폭을 만든다.
\[
\frac{\delta\rho}{\rho}\sim10^{-5},
\]
상세한 CMB 스펙트럼을 이용하면 소수의 우주론 매개변수를 여러 각도 규모에 걸쳐 검증할 수 있다.

\section{원시 씨앗에서 다시 하늘로}

이제 과거 빛 껍질들로 둘러싸인 중심의 관측자로 돌아가 보자. 마지막 산란 구면은 삼차원 원시 밀도장을 가로지른다. 그 구면이 서로 다른 중력 퍼텐셜을 지나는 방향에서는 마이크로파 온도가 조금씩 달라진다. WMAP은 그 구면 위의 이차원 무늬를 보여 준다. 이 지도는 시간의 시작 순간을 찍은 사진도 아니고, 구면 안에 있는 모든 구조를 투영한 것도 아니다. 마이크로파 광자가 마지막으로 산란된 바로 그 껍질 위의 온도 분포를 보여 준다.

껍질 안에서는 같은 원시 섭동들이 계속 진화한다. 중력이 과밀 영역을 증폭하고, 카우스틱과 비슷한 흐름이 시트와 필라멘트를 뚜렷하게 만들며, 기체가 모이고 은하가 형성된다. 따라서 CMB 무늬와 이후의 우주 거미줄은 하나의 역사가 거치는 두 단계다. 얼어붙은 인플레이션 요동은 먼저 작은 밀도 씨앗으로 나타나고, 나중에는 관측 가능한 구조가 된다.

\begin{classroomqa}
\audiencequestion{오늘날 우리가 보는 마이크로파 배경은 어느 시대를 나타내나요?}
\lecturerresponse{광자는 대략 \(13\)십억 년 동안 이동해 왔지만, 뜨거운 빅뱅이 일어난 뒤 수십만 년이 지나 전자와 양성자가 결합해 중성 수소를 이루고 우주가 투명해졌을 때 방출되었다. 인플레이션은 그보다 훨씬 먼저 끝났다. 재가열로 뜨거운 물질과 복사가 생겼고, 그 뒤 우주는 팽창하고 식어 마지막 산란에 이르렀다. 이 논의에서 확실히 말할 수 있는 것은 정확한 재결합 시점이 아니라 ---수십만 년--- 정도라는 규모다. \lecturetimestamp{02:00:38--02:02:19}}
\end{classroomqa}

\section{하나뿐인 하늘과 우주 분산}

\begin{classroomqa}
\audiencequestion{인플레이션 이론은 WMAP에서 보고된 선호 방향이나 남--북 비대칭을 예측하나요?}
\lecturerresponse{여기서 다룬 단순한 인플레이션 그림으로부터 그런 선호 방향이 따라 나오지는 않는다. 보고된 큰 각도 규모의 이상 현상은 그 지위가 불확실했고, 전경 신호와 통계적 요동도 여전히 고려해야 할 문제였으므로, 책임 있는 평가는 선정적인 요약이 아니라 전문 문헌을 살펴봐야 했다. 인플레이션 이론은 앙상블의 통계적 성질을 예측하지만, 하늘에서 특이해 보이는 현상 하나가 새로운 물리학의 증거인지 판단하기는 더 어렵다. \lecturetimestamp{02:02:19--02:03:26}}
\end{classroomqa}

작은 각도 규모의 특징은 하늘 곳곳에 여러 번 나타난다. 표본이 많으므로 그 분포와 등방성을 통계적으로 검증할 수 있다. 사중극자 규모의 특징은 다르다. 관측 가능한 구면에는 서로 독립적인 큰 각도 규모 모드가 몇 개밖에 들어가지 않는다. 따라서 근본 이론이 옳더라도 측정된 파워는 앙상블 평균에서 상당히 벗어날 수 있다. 이 피할 수 없는 표본 추출의 불확실성이 \emph{우주 분산}이다.

\begin{classroomqa}
\audiencequestion{낮은 다중극 모드 몇 개가 겉보기에 정렬되어 있는 것은 어떻게 봐야 하나요?}
\lecturerresponse{그것은 단서일 수도 있고, 전경 신호나 분석 과정에서 생긴 인공물일 수도 있으며, 가능하기는 하지만 드문 실현일 수도 있다. 어려운 점은 같은 통계적 앙상블에서 관측 가능한 우주를 하나 더 뽑을 수 없다는 것이다. 동전을 여섯 번 던져 네 번 앞면이 나왔다고 해서 동전이 치우쳤다고 단정할 수는 없다. 여러 차례 반복하면 판단할 수 있겠지만, 우주론에서는 가장 큰 규모에 대해 한 묶음의 표본밖에 얻지 못한다. 이 이상 현상은 연구할 가치가 있을 수 있지만, 하늘 하나만으로는 결정적인 유의성 검정에 필요한 반복 시행을 제공할 수 없다. \lecturetimestamp{02:05:31--02:07:36}}
\end{classroomqa}

\section{양자 운동에서 우주 구조까지}

세 가지 보편 상수로 플랑크 규모를 정의하면, 익숙한 입자 물리학이 양자 중력과 얼마나 멀리 떨어져 있는지 알 수 있다. 마이크로파 하늘은 정밀한 우주론적 목표를 제시한다. 거의 등방적인 하늘, \(10^{-5}\) 정도의 요동 진폭, 그리고 여러 규모에 걸쳐 퍼진 구조가 그것이다. 양자역학에 따르면 모든 장 모드에는 영점 요동이 있다. 인플레이션은 각 인플라톤 모드의 물리적 파장이 대략 \(H^{-1}\)가 될 때까지 늘린다. 이때 허블 감쇠가 그 진동을 억제하고, 사실상 동결된 진폭을 남긴다.

균일한 인플라톤은 동결된 요동 아래에서 계속 구른다. 장이 평탄부의 끝에 이르면 이웃한 영역들이 정확히 같은 시각에 재가열되지는 않는다. 먼저 끝난 영역은 더 오래 희석되고, 나중에 끝난 영역은 더 높은 밀도로 남는다. 이 과정은 \(\delta\phi\)를 원시 밀도 및 곡률 무늬로 바꾼다. 그러면 중력이 그 무늬를 증폭해 은하와 우주 거미줄을 만들고, 마지막 산란 껍질은 마이크로파 하늘에 그보다 이른 상태를 보존한다. 따라서 미시적인 영점 운동과 우주 팽창과 시차를 두고 일어나는 재가열은 서로 별개의 이야기가 아니다. 이들이 함께 인플레이션 우주가 어떻게 엄청나게 균일하면서도 우리가 존재할 만큼 충분한 구조를 갖게 되었는지 설명한다.
